
A. Critical Input Sensitivities (factores más críticos)

High-Demand Surplus Events
(Muchos usuarios intentan reclamar el mismo surplus al mismo tiempo → riesgo de race conditions y conflictos de asignación).
Large Volumes of Simultaneous Users
(Pico de concurrencia en horas pico → degradación de tiempos de respuesta y posible inconsistencia en la base de datos).
Geolocation Service Failure
(Falla externa de GPS o API de mapas → el algoritmo de matching se detiene porque no puede calcular distancias).
Confirmed Match with No Pickup
(Usuario acepta pero no recoge en el tiempo definido → surplus queda bloqueado y se pierde).
5. Ausencia de usuarios en el radio operativo o datos de ubicación obsoletos en el momento de aceptación
B. Emergent Behaviors (comportamientos que surgen con el uso)

Informal Priority Networks Among Users
(Grupos de WhatsApp o acuerdos externos que saltan el algoritmo y distorsionan la equidad).
Supplier Disengagement Under Low Demand
(Si los surpluses no se reclaman, los cafeterías/restaurantes dejan de publicar → ciclo negativo).
Inequitable Distribution Patterns
(El algoritmo favorece siempre a los usuarios más cercanos → usuarios periféricos se desconectan).

C. Chaotic System Dynamics (efectos no lineales y cascadas)

Nonlinear Cascade from a Single No-Show
(Un no-show genera re-asignaciones sucesivas que consumen todo el tiempo hasta que el surplus caduca).
Demand Amplification Through Notifications
(Notificación simultánea a muchos usuarios → pico de demanda artificial que colapsa el sistema).
Feedback Loop Between Metrics and Supplier Behavior
(El dashboard muestra bajas tasas → proveedores se desmotivan y publican menos → las métricas empeoran más).